\paragraph{\trAutomorphism。}\label{sub:automorphisms}
§10 中就 $\scr{F}(M)^{\mathrm{pr}}$ 所作的全部观察，都可承袭到 $\scr{F}(M)$ 上。当然，有些陈述须稍作改动以容纳\trBraneStructure，但这通常相当直截了当。作为一种实际演示，我们将考虑比较吸引人的例子之一，即\trSymplectic\trAutomorphism{}群在\trFukayaCategory{}上的作用——这最初在\cref{sub:autoequivalences}中讨论过。
我们先从一些一般性的讨论开始。取任意一个\trSymplectic\trManifold{}$(M,\omega)$，带有\trAlmostComplex{}结构 $I_M$ 与二次复\trVolume{}形式 $\eta_M^2$。$M$ 的一个\textbf{\trGraded\trAutomorphism}（graded automorphism）是一个二元组 $(\phi,\gamma^{\#})$，由一个\trSymplectic\trAutomorphism{}$\phi$ 与一个映射 $\gamma^{\#}\colon Gr(TM)\to\RR$ 组成，后者满足
\begin{equation*}
\exp(2\pi i\,\gamma^{\#}(\Lambda))=\frac{\alpha_M(D\phi(\Lambda))}{\alpha_M(\Lambda)}
\end{equation*}
对一切 $\Lambda\in Gr(TM)$。此类\trAutomorphism{}的\trComposition{}定义为 $\phi_2^{\#}\circ\phi_1^{\#}=(\phi_2\circ\phi_1,\ \gamma_2^{\#}\circ D\phi_1+\gamma_1^{\#})$。\trGraded\trAutomorphism{}自然地作用在由\trLagrange\trBrane{}组成的集合上：若 $(L,\alpha^{\#},P^{\#})$ 是一个膜，则像 $\phi(L)$ 承袭\trGrading{}$x\mapsto\gamma^{\#}(TL_{\phi^{-1}(x)})+\alpha^{\#}(\phi^{-1}(x))$，以及\trPushforward{} Pin~结构 $(\phi^{-1})^*P^{\#}$。这还算不上最一般的故事。定义一个\textbf{\trBrane\trAutomorphism}（brane automorphism）为三元组 $\phi^{\#}=(\phi,\gamma^{\#},\beta)$，其中 $(\phi,\gamma^{\#})$ 照前，而 $\beta\to M$ 是一个\trRealLineBundle。此类\trAutomorphism{}对膜的作用现在还包含 Pin~结构的一个扭转，
\begin{equation}\label{eq:brane-automorphism-pin-twist}
P^{\#}\mapsto(\phi^{-1})^*(P^{\#}\otimes\beta|L).
\end{equation}
一个与此相关的例子是\trSuspension{}操作 $S$：鉴于\cref{thm:grading-relative-orientation}，它可以写成 $(\mathrm{Id}_M,\ \gamma^{\#}=-1,\ \beta=\xi_M)$ 的作用。
从现在起，假设 $M$ 是一个带角的\trExactSymplectic\trManifold，并考虑由那些使底层 $\phi$ 落在 $\mathrm{Aut}(M,\partial M)$ 内的 $(\phi,\gamma^{\#})$ 组成的群 $\mathrm{Aut}^{\mathrm{gr}}(M,\partial M)$。我们赋予它由 $\mathrm{Aut}(M,\partial M)\times C^{\infty}(Gr(TM),\RR)$ 诱导的\trTopology。存在 $\mathrm{Aut}^{\mathrm{gr}}(M,\partial M)$ 在 $H\scr{F}(M)$ 上的一个\trCanonical{}严格作用，而且它下降为 $\pi_0(\mathrm{Aut}^{\mathrm{gr}}(M,\partial M))$ 的一个弱作用。对\trFukayaCategory{}本身而言，与\cref{sub:autoequivalences}及\cref{sub:hamiltonian-isotopies}中相同的构造给出 $\mathrm{Aut}^{\mathrm{gr}}(M,\partial M)$ 在 $\scr{F}(M)$ 上的一个\trCoherent{}作用，它同样下降为 $\pi_0(\mathrm{Aut}^{\mathrm{gr}}(M,\partial M))$ 的弱作用。现在把一般的\trBrane\trAutomorphism{}包括进来，但不区分 $M$ 上相互\trIsomorphic{}的\trLineBundle。这定义了 $\mathrm{Aut}^{\mathrm{gr}}(M,\partial M)$ 经由 $H^1(M;\ZZ/2)$ 的一个扩张，我们把它记作 $\mathrm{Aut}^{\#}(M,\partial M)$。该群在\trFukayaCategory、乃至\trDonaldson--\trFukayaCategory{}上的自然作用，一般而言将只是弱作用，原因无非是\eqref{eq:brane-automorphism-pin-twist} 需要在\trIsomorphism{}类 $[\beta]$ 之内选取一个代表 $\beta$。这一现象在实践中并不太要紧，但其痕迹，稍后当我们在 §14 讨论\trInvolution{}时，细心的眼睛仍会察觉。
\begin{remark}\label{thm:representative-choice-remark}
要看出代表 $\beta$ 的选取确实要紧，考虑如下情形：我们有三条\trExactLagrangianBrane{}$L_k^{\#}=(L_k,\alpha_k^{\#},P_k^{\#})$（$k=0,1,2$），其中后两条是\trIsomorphic{}的（享有同一个\trLagrange\trSubmanifold{}与同一个\trGrading，并带有\trIsomorphic{}但不一定严格相等的 Pin~结构）。Pin~\trIsomorphism{}$g\colon P_1^{\#}\to P_2^{\#}$ 的一个选取确定诸\trFloerCohomology{}群之间的一个\trIsomorphism
\begin{equation}\label{eq:pin-choice-floer-iso}
HF^*(L_0^{\#},L_1^{\#})\cong HF^*(L_0^{\#},L_2^{\#}).
\end{equation}
为简单计，假设 $L_0$ 与 $L_1=L_2$ \trTransverse{}相交，于是有关\trFloerDatum{}中的\trHamilton{}项可以取成零。取定一个交点 $y$，并令 $\Lambda_k^{\#}=\Lambda_{k,y}^{\#}$（$k=0,1,2$）为诸 $L_k^{\#}$ 在该点处给出的\trLinear\trBrane。在正确的\trHomotopy{}类中取一条路径 $\lambda\in\Omega^-(Gr(TM_x),\Lambda_0,\Lambda_1)$\footnote{译注：原书此处印作 $Gr(TM_x)$；疑应为 $Gr(TM_y)$（前文取定的交点为 $y$，诸 $\Lambda_k^{\#}$ 均为在该点处取得的\trLinear\trBrane）；此处照原书排印，存疑备案。——译者注}，并在关联丛 $F_\lambda\to[0;1]$ 上取一个 Pin~结构 $P_\lambda^{\#}$。可以把 $(\delta_H)_\lambda$ 同时与两个\trOrientation{}空间 $o(\Lambda_0^{\#},\Lambda_1^{\#})$ 与 $o(\Lambda_0^{\#},\Lambda_2^{\#})$ 等同起来，见\cref{thm:halfplane-index-determinant}。然而，这一等同依赖于一些辅助数据的选取，即\trIsomorphism{}$f_0\colon(P_\lambda^{\#})_0\to P_0^{\#}$ 与 $f_k\colon(P_\lambda^{\#})_1\to P_k^{\#}$（$k=1,2$）。要看出\eqref{eq:pin-choice-floer-iso} 底下的\trChainMap，应当取
\begin{equation*}
f_2=g_y\circ f_1
\end{equation*}
并考察由此诱导的\trOrientation{}空间之间的映射。把 $g$ 的符号翻转会影响 $f_2$ 的符号，从而（依照\cref{thm:brane-orientation-reversal}中的讨论）改变等同 $(\delta_H)_\lambda\cong o(\Lambda_0^{\#},\Lambda_2^{\#})$，最终把\eqref{eq:pin-choice-floer-iso} 乘上 $-1$。
改进这一情形的一个办法，是不把 $\mathrm{Aut}^{\#}(M,\partial M)$ 当作群、而是当作一个稍更精细的\trObject{}来考虑，即“弱 2-群”，它把诸\trLineBundle{}的内部对称群 $H^0(M;\ZZ/2)$ 也考虑在内。这同一般的\trCategory{}理论颇为合拍——后者本就会主张，“\trCategory{}的\trAutomorphism{}群”无论如何正是这样一种\trObject。用这种语言来说，\cref{thm:representative-choice-remark}中的计算表明：平凡\trLineBundle{}$\beta=M\times\RR$ 的\trAutomorphism{}$-\mathrm{Id}_\beta$ 给出从恒等\trFunctor{}到其自身的\trNaturalTransformation{}$-E$。
\end{remark}
